\subsection{রূপান্তরের কারণ, যাচাই ও সমন্বিত উদাহরণ}
ভাগের নিয়মটি আসে $N=bq+r$ থেকে। প্রথম ভাগশেষ এককের ঘরের অঙ্ক; পরের ভাগশেষ তার বামের অঙ্ক। এজন্য ভাগশেষ উল্টো ক্রমে পড়তে হয়। ভগ্নাংশে $bf=d+f_1$; এখানে $d$ প্রথম ভগ্নাংশ অঙ্ক, তাই গুণ থেকে পাওয়া পূর্ণাংশ সরাসরি ক্রমে পড়তে হয়।

\chThreeExample{১}{বাইনারি ভগ্নাংশ কেন বারবার চলতে পারে?}
$(0.1)_{10}$-কে বাইনারিতে লেখো।
\chThreeAnswer
\begin{center}\begin{tabular}{|c|c|c|}\hline
গুণ & প্রাপ্ত অঙ্ক & অবশিষ্ট ভগ্নাংশ\\\hline
$0.1\times2=0.2$ & 0 & 0.2\\\hline
$0.2\times2=0.4$ & 0 & 0.4\\\hline
$0.4\times2=0.8$ & 0 & 0.8\\\hline
$0.8\times2=1.6$ & 1 & 0.6\\\hline
$0.6\times2=1.2$ & 1 & 0.2\\\hline
\end{tabular}\end{center}
০.২ আবার এসেছে, তাই পরের অঙ্কগুলো পুনরাবৃত্ত হবে: $(0.1)_{10}=(0.000110011\ldots)_2$। প্রথম ৮ ভগ্নাংশ বিট নিলে $(0.00011001)_2=25/256=0.09765625$; এই সীমিত মানের আগে $\approx$ লিখতে হবে। সরলতম ভগ্নাংশের হরের মৌলিক গুণনীয়ক সবই লক্ষ্য ভিত্তির গুণনীয়ক হলে সেই ভিত্তিতে ভগ্নাংশটি সসীম হয়। $1/10$-এর হরে ৫ আছে, কিন্তু বাইনারির ভিত্তিতে নেই।

\chThreeExample{২}{একই সংখ্যা চার পদ্ধতিতে}
$(58.375)_{10}$-কে বাইনারি, অক্টাল ও হেক্সাডেসিমালে প্রকাশ করো।
\chThreeAnswer $58=32+16+8+2$, তাই পূর্ণাংশ $111010$। ভগ্নাংশে $0.375\times2=0.75$, $0.75\times2=1.5$, $0.5\times2=1$; অঙ্ক ০, ১, ১।
\[\begin{aligned}
(58.375)_{10}&=(111010.011)_2\\
&=(111\;010.011)_2=(72.3)_8\\
&=(0011\;1010.0110)_2=(3A.6)_{16}.
\end{aligned}\]
যাচাই: $3\times16+10+6/16=58.375$। grouping-এর জন্য বসানো শূন্য মান বদলায়নি।

\chThreeExample{৩}{নম্বর তুলনা ও গড়}
এক শিক্ষার্থীর তিন বিষয়ের নম্বর $(4C)_{16}$, $(115)_8$ ও $(1001000)_2$। সর্বোচ্চ নম্বর ও গড় নির্ণয় করো।
\chThreeAnswer
\[(4C)_{16}=4\times16+12=76,\quad(115)_8=64+8+5=77,\]
\[(1001000)_2=64+8=72.\]
সর্বোচ্চ নম্বর ৭৭; গড় $(76+77+72)/3=75=(1001011)_2$। তুলনা বা সাধারণ গড় করার আগে একই ভিত্তিতে আনা হয়েছে।

\chThreeExample{৪}{ভিন্ন ভিত্তির সংখ্যার ব্যবধান}
কলেজের দুটি শ্রেণিতে $(3B)_{16}$ ও $(67)_8$ জন শিক্ষার্থী আছে। ব্যবধান বাইনারিতে কত?
\chThreeAnswer $3\times16+11=59$, $6\times8+7=55$। ব্যবধান $4=(100)_2$। লিখিত অঙ্কের দৈর্ঘ্য দেখে সংখ্যার মান তুলনা করা ঠিক নয়।

\chThreeExample{৫}{স্থানমান যোগ করে দ্রুত রূপান্তর}
$(173)_{10}$-কে ভাগ না করে বাইনারিতে লেখো।
\chThreeAnswer $173=128+32+8+4+1$। তাই $128,64,32,16,8,4,2,1$-এর নিচে যথাক্রমে $1,0,1,0,1,1,0,1$ বসে।
\[(173)_{10}=(10101101)_2=(255)_8=(AD)_{16}.\]

\chThreeExample{৬}{ধারণাভিত্তিক MCQ: শূন্য কোথায় বসবে?}
$(101.11)_2$-এর হেক্সাডেসিমাল মান কোনটি? (ক) $(5.3)_{16}$, (খ) $(5.C)_{16}$, (গ) $(A.3)_{16}$, (ঘ) $(A.C)_{16}$।
\chThreeAnswer (খ)। grouping $0101.1100$। পূর্ণাংশে শূন্য বামে, ভগ্নাংশে ডানে; ভগ্নাংশের শুরুতে শূন্য বসালে মান বদলে যায়।

\chThreeExample{৭}{শুধু ভগ্নাংশের তিন রূপান্তর}
\chThreeAnswer
\begin{center}\begin{tabularx}{\linewidth}{|L{.23\linewidth}|Y|L{.27\linewidth}|}\hline
প্রদত্ত মান & ধারাবাহিক গুণ থেকে পাওয়া অঙ্ক & ফল\\\hline
$(0.8125)_{10}$ & $0.8125\times2=1.625$; $0.625\times2=1.25$; $0.25\times2=0.5$; $0.5\times2=1$ & $(0.1101)_2$\\\hline
$(0.046875)_{10}$ & $0.046875\times8=0.375$; $0.375\times8=3$ & $(0.03)_8$\\\hline
$(0.6875)_{10}$ & $0.6875\times16=11=B$ & $(0.B)_{16}$\\\hline
\end{tabularx}\end{center}
প্রথম উদাহরণে অঙ্ক ১, ১, ০, ১; দ্বিতীয়টিতে ০, ৩। ভগ্নাংশের শুরুর ০ বাদ দেওয়া যাবে না। তিন ক্ষেত্রেই পূর্ণাংশ শূন্য এবং ভগ্নাংশ সসীম।

% Mahbub teaching-type follow-up, 2026-10-14
\chThreeExample{৮}{শুধু বাইনারি, অক্টাল ও হেক্সাডেসিমাল ভগ্নাংশ থেকে দশমিক}
$(0.1011)_2$, $(0.47)_8$ ও $(0.A6)_{16}$-কে দশমিকে প্রকাশ করো।
\chThreeAnswer
\[\begin{aligned}
(0.1011)_2&=\tfrac12+\tfrac18+\tfrac1{16}=0.6875,\\
(0.47)_8&=\tfrac48+\tfrac7{64}=0.609375,\\
(0.A6)_{16}&=\tfrac{10}{16}+\tfrac6{256}=0.6484375.
\end{aligned}\]
বিন্দুর পর প্রথম অঙ্কের ওজন $b^{-1}$; $b^0$ নয়। পূর্ণাংশ না থাকলেও বিন্দু ও তার আগে ০ লেখা সুবিধাজনক।

\chThreeExample{৯}{শুধু ভগ্নাংশের grouping ও বিপরীত রূপান্তর}
একই ভগ্নাংশকে binary, octal ও hexadecimal-এ দেখাও।
\chThreeAnswer
\[\begin{aligned}
(0.10111)_2&=(0.101\;110)_2=(0.56)_8,\\
(0.10111)_2&=(0.1011\;1000)_2=(0.B8)_{16},\\
(0.56)_8&=(0.101\;110)_2=(0.10111)_2,\\
(0.B8)_{16}&=(0.1011\;1000)_2=(0.10111)_2.
\end{aligned}\]
অতএব octal থেকে hex এবং hex থেকে octal-এও যথাক্রমে $0.56_8\to0.101110_2\to0.B8_{16}$ ও বিপরীত পথ চলে। তিন পদ্ধতির মান $0.71875$। শেষের শূন্য বাদ দেওয়া যায়; ভগ্নাংশের \emph{শুরুর} শূন্য বাদ দিলে মান বদলায়।

\chThreeExample{১০}{স্থানমান বেছে অক্টাল ও হেক্সাডেসিমালে দ্রুত রূপান্তর}
ভাগ না করে $(345)_{10}$-কে octal ও hexadecimal-এ লেখো।
\chThreeAnswer অক্টালের ওজন ৬৪, ৮, ১: $345=5\times64+3\times8+1$; তাই $(531)_8$। হেক্সাডেসিমালের ওজন ২৫৬, ১৬, ১: $345=1\times256+5\times16+9$; তাই $(159)_{16}$। প্রতিটি নির্বাচিত অঙ্ক ভিত্তির চেয়ে ছোট থাকতে হবে। এটি একই স্থানমানের নিয়মের বিপরীত ব্যবহার।

\chThreeExample{১১}{ডিজিটাল সিগন্যাল থেকে সংখ্যামান}
আটটি সমান সময়খণ্ডে একটি সংকেতের level যথাক্রমে উচ্চ, নিম্ন, উচ্চ, উচ্চ, নিম্ন, নিম্ন, উচ্চ, নিম্ন। প্রথম খণ্ডকে MSB ধরে বাইনারি, অক্টাল ও হেক্সাডেসিমাল মান নির্ণয় করো।
\chThreeAnswer উচ্চ=১ ও নিম্ন=০ হলে bit sequence $10110010$।
\[(10110010)_2=(010\;110\;010)_2=(262)_8=(1011\;0010)_2=(B2)_{16}=178_{10}.\]
প্রথম খণ্ড LSB হলে পড়ার ক্রম উল্টে $01001101$ হবে; তখন মান ৭৭। তাই সিগন্যালের সময়ক্রম ও বিটের গুরুত্বের ক্রম আলাদা করে উল্লেখ করতে হয়।
